\section{Proofs for \texorpdfstring{Section~\ref{sec:erw}}{Section 2}}\label{app:erw}

Throughout this appendix $n\ge2$ is even, $h=n/2$, and we condition on
$X_1=+1$. So $M_1=0$, the chain $M$ moves from $b$ to $b+1$ at time $t$ with
probability $r_t(b)$, and $V_s(b)=\Prob_+(M_n=h\mid M_s=b)$. We write
$\mathcal L_sF(b)=(1-r_s(b))F(b)+r_s(b)F(b+1)$, so that
$V_s=\mathcal L_sV_{s+1}$. Both $M_t$ and $t-M_t\ge1$ are nondecreasing. So
from a state $b\ge1$ the chain stays at states $\ge1$, and for $b\ge1$ we
have $V_s(b)=0$ unless $b$ is feasible, that is,
$b\in\mathcal S_s=\{b\in\Z:1\le b\le h,\ 1\le s-b\le h\}$. Note that
$\mathcal S_n=\{h\}$. Theorem~\ref{thm:erw-bounds}(a) is
Proposition~\ref{prop:A-super}(d), and Theorem~\ref{thm:erw-bounds}(b) is
Proposition~\ref{prop:A-bridge}.

\subsection{Pseudo-count urns}\label{app:erw-urns}

For $c\ge0$ the P\'olya$(c)$ chain moves at time $s$ from $b$ to $b+1$ with
probability $p^c_s(b)=\frac{b+c}{s+2c}$ and stays at $b$ otherwise. It is
P\'olya's urn with weights $b+c$ and $s-b+c$ for the 2 colors, where each
draw adds a ball of the color drawn. For $b+c>0$ and $s-b+c>0$ let
\[
 \Phi^c_s(b)=\binom{n-s}{h-b}\frac{\mathrm B(h+c,h+c)}{\mathrm B(b+c,\,s-b+c)},
\]
where $\mathrm B$ is the beta function and the binomial coefficient is 0 unless
$0\le h-b\le n-s$. For $1\le b\le s-1$ this is defined for all $c\ge0$, and
it vanishes at infeasible $b$.

\begin{lemma}\label{lem:A-polya}
Let $c\ge0$, $1\le s\le n$, $b+c>0$ and $s-b+c>0$.
\begin{enumerate}[(a)]
\item $\Phi^c_s(b)$ is the probability that the P\'olya$(c)$ chain at $b$ at
time $s$ is at $h$ at time $n$. So $\Phi^c_n(b)=\mathbf 1\{b=h\}$, and
$\Phi^c_s(b)=p^c_s(b)\Phi^c_{s+1}(b+1)+(1-p^c_s(b))\Phi^c_{s+1}(b)$ for
$s<n$.
\item If $s<n$ and $\Phi^c_s(b)>0$, then
$p^c_s(b)\Phi^c_{s+1}(b+1)/\Phi^c_s(b)=r^{\br}_s(b):=\frac{h-b}{n-s}$. So,
given that it ends at $h$, the chain moves up with probability
$r^{\br}_s(b)$ for every $c$, and its remaining $h-b$ up steps occupy a
uniformly random set of $h-b$ of the remaining $n-s$ steps.
\item If $0<\varepsilon<\frac12$ and $\beta=\frac{\varepsilon}{1-2\varepsilon}$,
then $r_s(b)=p^{\beta s}_s(b)$ for $0\le b\le s$. So at time $s$ the chain
$M$ moves like the P\'olya$(\beta s)$ chain.
\item If $0<\varepsilon<\frac12$ and $0\le c\le\beta s$, then
$V_s(b)\ge\Phi^c_s(b)$ for $1\le b\le s-1$. In particular
$V_{t+1}(1)\ge\Phi^0_{t+1}(1)=\pi_t$ for $1\le t\le h$.
\end{enumerate}
\end{lemma}

\begin{proof}
(a) Let $x=b+c$, $y=s-b+c$ and $(x)_k=x(x+1)\cdots(x+k-1)$. Each sequence
of $n-s$ draws with $h-b$ draws of the first color has probability
$(x)_{h-b}(y)_{n-s-h+b}/(x+y)_{n-s}=B(h+c,h+c)/B(x,y)$, and there are
$\binom{n-s}{h-b}$ such sequences. The recursion is the Markov property.

(b) Since $\binom{n-s-1}{h-b-1}/\binom{n-s}{h-b}=\frac{h-b}{n-s}$ and
$\mathrm B(x,y)/\mathrm B(x+1,y)=\frac{x+y}x$, we get
$\Phi^c_{s+1}(b+1)/\Phi^c_s(b)=\frac{h-b}{n-s}\cdot\frac{s+2c}{b+c}$.
Choosing a uniform set of $h-b$ of the $n-s$ steps one step at a time gives
the same up-probabilities.

(c) Since $1+2\beta=1/(1-2\varepsilon)$, we have
$\frac{b+\beta s}{s+2\beta s}=(1-2\varepsilon)(\frac bs+\beta)=r_s(b)$.

(d) Let $F_u=\Phi^c_u$ on $\{1,\dots,u-1\}$ for $s\le u\le n$. By (a),
\[
 \mathcal L_uF_{u+1}(b)-F_u(b)=\bigl(r_u(b)-p^c_u(b)\bigr)\bigl(F_{u+1}(b+1)-F_{u+1}(b)\bigr).
\]
By (c), $r_u(b)-p^c_u(b)=(\frac u2-b)(\frac1{u+2c}-\frac1{u(1+2\beta)})$,
and the last factor is $\ge0$ since $c\le\beta u$. Let $b\in\mathcal S_u$.
If $b=h$ then $b\ge u/2$ and $F_{u+1}(b+1)=0$. If $u-b=h$ then $b\le u/2$
and $F_{u+1}(b)=0$. Otherwise
$F_{u+1}(b+1)/F_{u+1}(b)=\frac{(h-b)(u-b+c)}{(h-u+b)(b+c)}$, which is $\ge1$
for $b\le u/2$ and $\le1$ for $b\ge u/2$. In each case the 2 factors above
have the same sign or one of them is 0, so $\mathcal L_uF_{u+1}\ge F_u$ on
$\mathcal S_u$. Since $F_n(b)=\mathbf 1\{b=h\}$ and $F_u$ vanishes off
$\mathcal S_u$, Lemma~\ref{lem:A-compare}(b) below gives $V_u\ge F_u$.
Finally $\Phi^0_{t+1}(1)=t\binom{n-t-1}{h-1}(h-1)!^2/(n-1)!=\pi_t$.
\end{proof}

\begin{lemma}\label{lem:A-compare}
Let $1\le s_0\le n$, and for $s_0\le s\le n$ let $F_s\ge0$ be a function on
$\{1,\dots,s-1\}$.
\begin{enumerate}[(a)]
\item If $F_n(h)\ge1$ and $\mathcal L_sF_{s+1}\le K_sF_s$ on $\mathcal S_s$
for $s_0\le s<n$, with constants $K_s\ge0$, then
$V_s\le\bigl(\prod_{u=s}^{n-1}K_u\bigr)F_s$ on $\mathcal S_s$ for
$s_0\le s\le n$.
\item If $F_n(b)\le\mathbf 1\{b=h\}$, each $F_s$ vanishes off $\mathcal S_s$,
and $\mathcal L_sF_{s+1}\ge F_s$ on $\mathcal S_s$ for $s_0\le s<n$, then
$V_s\ge F_s$ for $s_0\le s\le n$.
\end{enumerate}
\end{lemma}

\begin{proof}
We use backward induction on $s$, starting from
$V_n(b)=\mathbf 1\{b=h\}$ and using $V_s=\mathcal L_sV_{s+1}$ on
$\mathcal S_s$. In (a), $V_{s+1}$ vanishes at infeasible states and
$F_{s+1}\ge0$, so the induction hypothesis gives
$V_s\le\prod_{u>s}K_u\,\mathcal L_sF_{s+1}\le\prod_{u\ge s}K_u\,F_s$ on
$\mathcal S_s$. In (b) it gives $V_s\ge\mathcal L_sF_{s+1}\ge F_s$ on
$\mathcal S_s$, and $F_s=0\le V_s$ off $\mathcal S_s$.
\end{proof}

\subsection{The supersolution}\label{app:erw-super}

Let $\psi$ be the digamma function.

\begin{lemma}\label{lem:A-cderiv}
Let $2\le s\le n$, $b\in\mathcal S_s$ and $\tau\ge0$. Put $X=b+\tau$,
$Y=s-b+\tau$, $\Sigma=X+Y$ and $z=(s-2b)^2/\Sigma^2$. Then
$\partial_\tau\log\Phi^\tau_s(b)=2\psi(\Sigma)-\psi(X)-\psi(Y)
-2\psi(n+2\tau)+2\psi(h+\tau)$, and
\[
 \partial_\tau\log\Phi^\tau_s(b)
 <-\log(1-z)+\frac{1+z}{(1-z)\Sigma}+\frac{2(1+z)}{3\Sigma^2(1-z)^2}.
\]
\end{lemma}

\begin{proof}
The formula follows from $\log\mathrm B(x,y)=\log\Gamma(x)+\log\Gamma(y)
-\log\Gamma(x+y)$. For $x>0$, \cite[5.9.15]{dlmf} gives
$\psi(x)=\log x-\frac1{2x}-2\int_0^\infty\frac{t\,dt}{(t^2+x^2)(e^{2\pi t}-1)}$,
and the subtracted term is positive and less than
$\frac2{x^2}\int_0^\infty\frac{t\,dt}{e^{2\pi t}-1}=\frac1{12x^2}$. So
$\log x-\frac1{2x}-\frac1{12x^2}<\psi(x)<\log x-\frac1{2x}$. We use the upper
bound for $\psi(\Sigma)$ and the lower bound for $\psi(X)$ and $\psi(Y)$. The
logarithmic derivative of the duplication formula \cite[5.5.5]{dlmf} is
$2\psi(2y)=\psi(y)+\psi(y+\frac12)+2\log2$, and $\psi$ is increasing, so
$2\psi(n+2\tau)-2\psi(h+\tau)>2\log2$. Hence the derivative is less than
$\log\frac{\Sigma^2}{4XY}+(\frac1{2X}+\frac1{2Y}-\frac1\Sigma)
+\frac1{12}(\frac1{X^2}+\frac1{Y^2})$. Now $4XY=\Sigma^2(1-z)$ and
$X^2+Y^2=\Sigma^2(1+z)/2$, so the 3 terms are $-\log(1-z)$,
$\frac{X^2+Y^2}{2XY\Sigma}=\frac{1+z}{(1-z)\Sigma}$ and
$\frac{X^2+Y^2}{12X^2Y^2}=\frac{2(1+z)}{3\Sigma^2(1-z)^2}$.
\end{proof}

\begin{proposition}\label{prop:A-super}
Let $0<\varepsilon<\frac14$, $\kappa=\frac{2\varepsilon}{1-4\varepsilon}$,
$c_s=\kappa(s+10)$ and $\Sigma_s=s+2c_s$. For $1\le s\le n-1$ and
$b\in\mathcal S_s$ let
$D_s(b)=\mathcal L_s\Phi^{c_{s+1}}_{s+1}(b)\big/\Phi^{c_s}_s(b)$.
\begin{enumerate}[(a)]
\item Let $c=c_s$, $c'=c_{s+1}$, $d=\frac s2-b$, $u=\frac{d}{s+2c'}$,
$v=\frac d{n-s}$ and $\rho=\frac{s+2c'}{s(1+2\beta)}$. Then $|u|<\frac12$
and
\[
 D_s(b)=\mathcal J\,\frac{\Phi^{c'}_s(b)}{\Phi^{c}_s(b)},\qquad
 \mathcal J=1-(\rho-1)\frac{u(u+v)}{\frac14-u^2}.
\]
\item $\log D_s(b)\le\frac{\kappa}{\Sigma_s}+\frac{2\kappa}{3\Sigma_s^2}$.
\item For $1\le t\le h$,
$V_{t+1}(1)\le\pi_t\exp\bigl(\kappa\bigl[H_{n-1}-H_t+\frac2{3t}\bigr]
+c_{t+1}\bigl[\log\frac{(t+1)^2}{4t}+\frac23\bigr]\bigr)$.
\item $C_n\le U_n$, with $U_n$ as in Theorem~\ref{thm:erw-bounds}(a).
\end{enumerate}
\end{proposition}

The slope $\kappa$ is chosen so that $2(\kappa-\beta)/(1+2\beta)=\kappa$,
which follows from the formulas for $\kappa$ and $\beta$. In the proof of (b)
this makes the gain in $\mathcal J$ cancel the main loss from raising the
pseudo-count from $c$ to $c'$, and the offset 10 in $c_s$ controls what is
left.

\begin{proof}
(a) Write $p'=p^{c'}_s(b)$, $\bar r=r^{\br}_s(b)$ and $r=r_s(b)$. By
Lemma~\ref{lem:A-polya}(a) and (b) with $c'$, the numerator of $D_s(b)$
divided by $\Phi^{c'}_s(b)$ is
$\frac{(1-r)(1-\bar r)}{1-p'}+\frac{r\bar r}{p'}$. Here $p'=\frac12-u$,
$\bar r=\frac12+v$ and $r=\frac12-\rho u$, since
$\frac12-r=(1-2\varepsilon)\frac ds=\rho u$. Over the denominator
$p'(1-p')=\frac14-u^2$ the numerator is $f(u,v)+f(-u,-v)$ with
$f(u,v)=(\frac12+\rho u)(\frac12-v)(\frac12-u)$, and this equals
$\frac14-\rho u^2-(\rho-1)uv$. This gives $\mathcal J$. Also
$|d|\le\frac s2-1$, so $|u|<\frac12$.

(b) Write $c=c_s$, $c'=c+\kappa$ and $\Sigma=\Sigma_s$. Since
$c'-\beta s=(\kappa-\beta)s+11\kappa$, the identity for $\kappa$ gives
$\rho-1=\kappa+\nu_s$ with $\nu_s=\frac{22\kappa}{s(1+2\beta)}$. As $u$ and
$v$ have the sign of $d$, $u(u+v)\ge u^2$, so
$\mathcal J\le1-(\kappa+\nu_s)W'$ with $W'=\frac{z'}{1-z'}$ and
$z'=4u^2=\frac{(s-2b)^2}{(s+2c')^2}$. Also $\mathcal J>0$ by the expression
in (a), since $0<r<1$ and $\bar r$, $1-\bar r$ are not both 0. So
$\log\mathcal J\le-(\kappa+\nu_s)W'$. Next,
$\log(\Phi^{c'}_s(b)/\Phi^c_s(b))$ is the integral over $\tau\in[c,c']$ of
the derivative in Lemma~\ref{lem:A-cderiv}. There $z$ decreases and $\Sigma$
increases in $\tau$, and the bound of the lemma increases in $z$ and
decreases in $\Sigma$, so we may take it at $\tau=c$. Let $X$, $Y$ and $z$
be as in the lemma at $\tau=c$, so that $X+Y=\Sigma$, and put
$W=\frac z{1-z}$. Then $\frac1{1-z}=1+W$, $\frac{1+z}{1-z}=1+2W$ and
$-\log(1-z)\le W$, so
\begin{equation}\label{eq:A-logD}
 \log D_s(b)\le\kappa\Bigl[W+\frac{1+2W}\Sigma+\frac{2(1+2W)(1+W)}{3\Sigma^2}\Bigr]
 -(\kappa+\nu_s)W'.
\end{equation}
We use 3 facts. (i) $X,Y\ge1+c$, since $b\ge1$ and $s-b\ge1$. So
$XY\ge(1+c)\Sigma/2$ and $1+W=\frac{\Sigma^2}{4XY}\le\frac\Sigma{2(1+c)}$.
(ii) $z'=\lambda z$ with $\lambda=(\frac\Sigma{\Sigma+2\kappa})^2
\ge1-\frac{4\kappa}\Sigma$. So $W-W'=\frac{(1-\lambda)z}{(1-z)(1-\lambda z)}
\le(1-\lambda)W(1+W)\le\frac{2\kappa}{1+c}W$ by (i). As $c\ge11\kappa$, this
gives $W'\ge\frac9{11}W$. (iii) $\Sigma\ge s$ and $1+c>\kappa s$. If $W>0$
then $s\ne2b$, which together with $b\ge1$ and $s-b\ge1$ forces $s\ge3$, so
$\Sigma^2\ge3s$.

By (ii) and (iii), $\kappa(W-W')\le\frac{2\kappa^2}{1+c}W\le\frac{2\kappa}sW$
and $\frac{2\kappa W}\Sigma\le\frac{2\kappa}sW$. Since
$(1+2W)(1+W)=1+3W+2W^2$, $\kappa$ times the last term in the bracket of
\eqref{eq:A-logD} equals
$\frac{2\kappa}{3\Sigma^2}+\frac{2\kappa W}{\Sigma^2}
+\frac{4\kappa W^2}{3\Sigma^2}$. Here
$\frac{2\kappa W}{\Sigma^2}\le\frac{2\kappa}{3s}W$ by (iii), and
$\frac{4\kappa W^2}{3\Sigma^2}\le\frac{4\kappa W(1+W)}{3\Sigma^2}
\le\frac{2\kappa W}{3\Sigma(1+c)}\le\frac{2\kappa}{3s}W$ by (i). With
$-\nu_sW'\le-\frac9{11}\nu_sW$ this gives
\[
 \log D_s(b)\le\frac\kappa\Sigma+\frac{2\kappa}{3\Sigma^2}
 +W\Bigl[\frac{16\kappa}{3s}-\frac{18\kappa}{s(1+2\beta)}\Bigr],
\]
and the bracket is negative because $1+2\beta<2$ for $\varepsilon<\frac14$.

(c) Apply Lemma~\ref{lem:A-compare}(a) with $s_0=t+1$, $F_s=\Phi^{c_s}_s$
and $K_s=\exp(\frac\kappa s+\frac{2\kappa}{3s^2})$, which is allowed by (b)
and $\Sigma_s\ge s$. Since $1\in\mathcal S_{t+1}$ and
$\sum_{s>t}s^{-2}\le\frac1t$, we get
$V_{t+1}(1)\le e^{\kappa[H_{n-1}-H_t+2/(3t)]}\,\Phi^{c_{t+1}}_{t+1}(1)$. As
$\pi_t=\Phi^0_{t+1}(1)$, the logarithm of $\Phi^{c_{t+1}}_{t+1}(1)/\pi_t$ is
the integral over $\tau\in[0,c_{t+1}]$ of the derivative in
Lemma~\ref{lem:A-cderiv}, whose bound is largest at $\tau=0$. There $X=1$,
$Y=t$ and $\frac1{1-z}=\frac{(t+1)^2}{4t}$, and the bound equals
$\log\frac{(t+1)^2}{4t}+\frac{t^2+1}{2t(t+1)}+\frac{t^2+1}{12t^2}
\le\log\frac{(t+1)^2}{4t}+\frac23$.

(d) Insert (c) into Proposition~\ref{prop:first-mutation}(a). Since
$c_{t+1}=\kappa(t+11)$, the exponent in (c) is
$\kappa H_{n-1}+\kappa\mathcal G(t)$, and $\pi_t=c_1w_t$.
\end{proof}

\subsection{The bridge formula}\label{app:erw-bridge}

By Lemma~\ref{lem:A-polya}(b), the law of the path under $\E^{\br}_t$ is
the P\'olya$(c)$ chain from $1$ at time $t+1$ conditioned on $M_n=h$, for
every $c$. At time $u$ it moves up with probability $r^{\br}_u(M_u)$. We
write $\Prob^{\br}_t$ for the matching probability. For $1\le b\le u-1$ let
\[
 \varpi^+_u(b)=1+\varepsilon\frac{u-2b}{b},\qquad
 \varpi^-_u(b)=1-\varepsilon\frac{u-2b}{u-b},\qquad
 Q_u(b)=(u-2b)^2\Bigl[\frac{r^{\br}_u(b)}{b^2}+\frac{1-r^{\br}_u(b)}{(u-b)^2}\Bigr],
\]
and let $\varpi_u=\varpi^+_u(M_u)$ if $M_{u+1}=M_u+1$ and
$\varpi_u=\varpi^-_u(M_u)$ otherwise.

\begin{proposition}\label{prop:A-bridge}
Let $0\le\varepsilon\le1$ and $1\le t\le h$. Then
$V_{t+1}(1)=\pi_t\,\E^{\br}_t\bigl[\prod_{u=t+1}^{n-1}\varpi_u\bigr]$. If
$0<\varepsilon\le\frac12$, then
$\log V_{t+1}(1)\ge\log\pi_t+\varepsilon A_t-\varepsilon^2B_t$, where
\begin{equation}\label{eq:A-Bt}
 B_t=\sum_{u=t+1}^{n-1}\E^{\br}_t\,Q_u(M_u)\ \ge0 ,
\end{equation}
and hence $C_n\ge J_n$, where
\begin{equation}\label{eq:A-Jn}
 J_n=\varepsilon\sum_{t=1}^h(1-\varepsilon)^{t-1}\pi_t
 \exp\bigl(\varepsilon A_t-\varepsilon^2B_t\bigr).
\end{equation}
\end{proposition}

\begin{proof}
For $u\ge t+1$ and $b\in\mathcal S_u$ put $R_u(b)=V_u(b)/\Phi^0_u(b)$. By
Lemma~\ref{lem:A-polya}(a) and (b) with $c=0$,
$\frac bu\Phi^0_{u+1}(b+1)=r^{\br}_u(b)\Phi^0_u(b)$ and
$(1-\frac bu)\Phi^0_{u+1}(b)=(1-r^{\br}_u(b))\Phi^0_u(b)$. Dividing
$V_u=\mathcal L_uV_{u+1}$ by $\Phi^0_u(b)$ gives
\[
 R_u(b)=(1-r^{\br}_u(b))\,\frac{1-r_u(b)}{1-b/u}\,R_{u+1}(b)
 +r^{\br}_u(b)\,\frac{r_u(b)}{b/u}\,R_{u+1}(b+1),
\]
where a term with an infeasible state has bridge probability 0. Also
$\frac{r_u(b)}{b/u}=\varpi^+_u(b)$ and
$\frac{1-r_u(b)}{1-b/u}=\varpi^-_u(b)$. So
$R_u(b)=\E^{\br}_t[\varpi_uR_{u+1}(M_{u+1})\mid M_u=b]$, and $R_n(h)=1$.
Iterating from $u=t+1$ and using $\Phi^0_{t+1}(1)=\pi_t$ gives the formula.

Now let $0<\varepsilon\le\frac12$. If $b\le u/2$ then $\varpi^+_u(b)\ge1$
and $\varpi^-_u(b)-1\in[-\varepsilon,0]$, and if $b\ge u/2$ the same holds
with $+$ and $-$ exchanged. So $\varpi_u-1\ge-\frac12$. By Jensen's
inequality, $\E\prod_u\varpi_u\ge\exp\sum_u\E\log\varpi_u$. Also
$\log(1+y)\ge y-y^2$ for $y\ge-\frac12$, since the difference vanishes at 0
and its derivative $y(1+2y)/(1+y)$ has the sign of $y$. With
$\bar r=r^{\br}_u(b)$ and $u\bar r-b=\frac{n(u-2b)}{2(n-u)}$,
\[
 \E^{\br}_t[\varpi_u-1\mid M_u=b]=\varepsilon(u-2b)\frac{u\bar r-b}{b(u-b)}
 =\varepsilon\theta_u(b),\qquad
 \E^{\br}_t[(\varpi_u-1)^2\mid M_u=b]=\varepsilon^2Q_u(b).
\]
Summing over $u$ gives the bound, and
Proposition~\ref{prop:first-mutation}(a) gives $C_n\ge J_n$.
\end{proof}

Under $\Prob^{\br}_t$ let $Z=M_u-1$ and $Z'=m-Z$ be the numbers of up steps
and of other steps among the $m=u-t-1$ steps between times $t+1$ and $u$. So
$Z$ is hypergeometric, with $m$ draws from $n-t-1$ steps of which $h-1$ are
up steps. Let $\Omega_u=\{u/8\le M_u\le7u/8\}$.

\begin{lemma}\label{lem:A-AB}
Let $n\ge16$, $1\le t\le h$ and $t+1\le u\le n-1$.
\begin{enumerate}[(a)]
\item $\displaystyle\E^{\br}_t(u-2M_u)^2=\frac{(t-1)^2(n-u)^2}{(n-t-1)^2}
+\frac{4(u-t-1)(h-1)(h-t)(n-u)}{(n-t-1)^2(n-t-2)}\le(t-1)^2+u$.
\item If $t\le\frac{n+2}4$, then
$A_t\ge2(H_{n-1}-H_t-1-\frac1t)\ge2(\log n-\log t-3)$.
\item If $t\le\frac n8$ and $u\ge8(t+1)$, then
$\Prob^{\br}_t(\Omega_u^c)\le2e^{-0.0943u}$.
\item If $t\le\frac n8$, then $B_t\le84(t+1+\log n)+683(t+1)^3+19100$.
\end{enumerate}
\end{lemma}

\begin{proof}
(a) Since $n-2=2(h-1)$, the mean $\E Z=\frac{m(h-1)}{n-t-1}$ gives
$\E(u-2M_u)=\frac{(t-1)(n-u)}{n-t-1}$, and the second term is
$4\Var Z$. For the bound use $n-u\le n-t-1$,
$4(h-1)(h-t)=(n-t-1)^2-(t-1)^2$, and $m(n-u)\le u(n-t-2)$ for $u\ge t+2$.

(b) Since $b(u-b)\le u^2/4$, $\theta_u(b)\ge\frac{2n(u-2b)^2}{(n-u)u^2}$.
Keeping only $4\Var Z$ in (a) gives
$A_t\ge P_t\sum_{u=t+1}^{n-1}\frac{u-t-1}{u^2}$ with
$P_t=\frac{2n(n-2)(n-2t)}{(n-t-1)^2(n-t-2)}$. Here $P_1=\frac{2n}{n-3}>2$,
and for real $t\in[1,\frac{n+2}4]$,
$\frac{d}{dt}\log P_t=-\frac2{n-2t}+\frac2{n-t-1}+\frac1{n-t-2}
\ge\frac3{n-t-1}-\frac2{n-2t}\ge0$, since the last step is equivalent to
$4t\le n+2$. So $P_t\ge2$. The sum is nonnegative and at least
$H_{n-1}-H_t-(t+1)\sum_{u>t}u^{-2}\ge H_{n-1}-H_t-\frac{t+1}t$. For the
second inequality use $H_{n-1}-H_t\ge\log\frac n{t+1}\ge\log n-\log t-\frac1t$.

(c) Here $m\ge\frac78u$, so $\frac u8\le\frac m7$. Since $t\ge1$,
$\E Z\ge\frac m2$, and since $t\le\frac n8$,
$\E Z'=\frac{m(h-t)}{n-t-1}\ge\frac{3m}8$. Hoeffding's inequality also holds
for sampling without replacement \cite[Section~6]{hoeffding1963}, so
$\Prob(Z-\E Z\le-\sigma m)\le e^{-2\sigma^2m}$, and the same holds for $Z'$.
Now $\{M_u<\frac u8\}\subset\{Z-\E Z<-\frac{5m}{14}\}$, and since
$u-M_u=t+Z'$, also $\{u-M_u<\frac u8\}\subset\{Z'-\E Z'<-\frac{13m}{56}\}$.
So $\Prob(\Omega_u^c)\le e^{-2(5/14)^2m}+e^{-2(13/56)^2m}\le2e^{-0.10778m}
\le2e^{-0.0943u}$.

(d) Since $\bar r(u-b)^2+(1-\bar r)b^2\le u^2$, we have
$Q_u(b)\le\frac{(u-2b)^2u^2}{b^2(u-b)^2}$. On $\Omega_u$,
$b(u-b)\ge\frac{7u^2}{64}$, so $Q_u\le\frac{4096}{49}\frac{(u-2b)^2}{u^2}
\le83.6\frac{(u-2b)^2}{u^2}$,
and by (a) these parts of $B_t$ add up to at most
$83.6\sum_{u>t}(\frac{(t-1)^2}{u^2}+\frac1u)\le83.6(t+H_{n-1})
\le84(t+1+\log n)$. Off $\Omega_u$, $b(u-b)\ge u-1\ge\frac u2$ gives
$Q_u\le4u^2$. For $u<8(t+1)$ these parts add up to at most
$\frac43(8(t+1))^3\le683(t+1)^3$. For $u\ge8(t+1)$, (c) bounds them by
$8u^2e^{-0.0943u}$, and $8\sum_{u\ge1}u^2e^{-0.0943u}<19100$.
\end{proof}

\subsection{Proof of \texorpdfstring{Theorem~\ref{thm:erw-bounds}(c)}{Theorem 2.6(c)}}\label{app:erw-proof-c}

Let $n\ge64$ be even and $0<\varepsilon\le\frac1{100(1+\log n)}$. Then
$\varepsilon<0.002$, so $\kappa\le2.084\,\varepsilon$ and
$\kappa\log(h+1)\le2.084\,\varepsilon\log n<0.021<\frac{\log2}{24}$. Since
$\pi_t=\frac t{n-t}\prod_{i=1}^{t-1}\frac{h-i}{n-i}$ and
$\frac{h-i}{n-i}\le\frac12$,
\begin{equation}\label{eq:A-wt}
 w_t=\tfrac{n+2}4\,\pi_t\le\tfrac{n+2}{2(n-t)}\,t2^{-t}\le1.032\,t2^{-t}
 \qquad(1\le t\le h).
\end{equation}

\emph{Upper bound.} By Theorem~\ref{thm:erw-bounds}(a),
$(1-\varepsilon)^{t-1}\le1$ and $\log(1+y)\le y$,
$\log C_n\le\log(\varepsilon c_1)+\kappa H_{n-1}
+\sum_tw_t(e^{\kappa\mathcal G(t)}-1)$. Since
$\log\frac{(t+1)^2}{4t}\ge0$ and $\frac23(t+11)>t\ge H_t$, we have
$\mathcal G(t)>0$. Since $\frac{(t+1)^2}{4t}\le\frac{t+1}2$,
$\frac23<\log2$ and $H_t\ge\frac2{3t}$, also
$\mathcal G(t)\le(t+11)\log(t+1)\le12t\log(t+1)$. So
$0<\kappa\mathcal G(t)\le\frac t2\log2$ for $t\le h$, and
$e^{\kappa\mathcal G(t)}-1\le\kappa\mathcal G(t)2^{t/2}$. By \eqref{eq:A-wt},
\[
 \sum_tw_t\bigl(e^{\kappa\mathcal G(t)}-1\bigr)
 \le1.032\cdot2.084\,\varepsilon\sum_{t\ge1}t(t+11)\log(t+1)2^{-t/2}
 \le563.3\,\varepsilon,
\]
since the last sum is $261.9013\ldots\le261.91$. Also
$\kappa\le2\varepsilon(1+8\varepsilon)$ and $H_{n-1}\le1+\log n$, so
$\kappa H_{n-1}\le2\varepsilon\log n+2\varepsilon+16\varepsilon^2(1+\log n)
\le2\varepsilon\log n+2.2\,\varepsilon$. As $c_1=\frac4{n+2}$, the upper
bound holds with $565.5\,\varepsilon\le566\,\varepsilon$.

\emph{Lower bound.} Let $T=\lfloor n/8\rfloor\ge8$. By
Proposition~\ref{prop:first-mutation}(a), Proposition~\ref{prop:A-bridge}
for $t\le T$, Lemma~\ref{lem:A-polya}(d) for $t>T$ and Jensen's inequality
for the probability vector $(w_t)$,
\[
 \log C_n\ge\log(\varepsilon c_1)+\sum_tw_t(t-1)\log(1-\varepsilon)
 +\sum_{t\le T}w_t\bigl(\varepsilon A_t-\varepsilon^2B_t\bigr).
\]
By the inequality $\log(1+y)\ge y-y^2$ used above and
Proposition~\ref{prop:first-mutation}(c), the first sum is at least
$-(\varepsilon+\varepsilon^2)\sum_tw_t(t-1)\ge-2\varepsilon-2\varepsilon^2$.
Next $T\le\frac{n+2}4$, so $A_t\ge2(\log n-\log t-3)$ for $t\le T$ by
Lemma~\ref{lem:A-AB}(b). Let $w_{\le T}=\sum_{t\le T}w_t$. By Jensen's
inequality $\sum_tw_t\log t\le\log\sum_tw_tt\le\log3$, so
\[
 \sum_{t\le T}w_t\,\varepsilon A_t\ge2\varepsilon(\log n-3)w_{\le T}-2\varepsilon\log3
 \ge2\varepsilon\log n-2\varepsilon(1-w_{\le T})\log n-6\varepsilon-2\varepsilon\log3 .
\]
By \eqref{eq:A-wt}, $1-w_{\le T}\le1.032\sum_{t>T}t2^{-t}=1.032(T+2)2^{-T}$,
and $n<8(T+1)$. As $(T+2)2^{-T}\log(8T+8)$ decreases for $T\ge8$,
$2\varepsilon(1-w_{\le T})\log n\le2.064\cdot10\cdot2^{-8}\log72\,
\varepsilon<0.35\,\varepsilon$. Finally Lemma~\ref{lem:A-AB}(d) and
Proposition~\ref{prop:first-mutation}(c) give
$\sum_{t\le T}w_tB_t\le84(3+1+\log n)+683\cdot124+19100
=84\log n+104128$. Collecting terms,
\[
 \log C_n-\log\frac{4\varepsilon}{n+2}-2\varepsilon\log n
 \ge-(2+0.35+6+2\log3)\,\varepsilon-\varepsilon^2(84\log n+104130),
\]
which gives the lower bound since $2+0.35+6+2\log3<11$. For the last claim,
the hypothesis on $\varepsilon$ gives
$\varepsilon^2(84\log n+1.05\cdot10^5)\le\varepsilon\,
\frac{84\log n+1.05\cdot10^5}{100(1+\log n)}$. The fraction decreases in
$n$, and its value at $n=64$ is $204.2\ldots$. Since $11+205<566$, the
two-sided bound follows. \qed

\subsection{Proof of \texorpdfstring{Theorem~\ref{thm:K2}}{Theorem 2.9}}\label{app:erw-K2}

\emph{Part (a).} By Proposition~\ref{prop:first-mutation}(a) and (b),
$c_2=\sum_t[-(t-1)\pi_t+\partial_\varepsilon V_{t+1}(1)|_{\varepsilon=0}]$.
In Proposition~\ref{prop:A-bridge} the bridge law does not depend on
$\varepsilon$, and each $\varpi_u$ is affine in $\varepsilon$ and equals 1 at
$\varepsilon=0$. So by the conditional mean computed in its proof,
$\partial_\varepsilon V_{t+1}(1)|_{\varepsilon=0}
=\pi_t\sum_u\E^{\br}_t[(\varpi_u-1)/\varepsilon]=\pi_tA_t$. Divide by $c_1$.

\emph{Part (b), the limit.} Fix $t$ and let $n\ge16(t+1)$. For
$1\le b\le u-1$,
\[
 \theta_u(b)=\frac n{n-u}\Bigl[\frac{2(u-2b)^2}{u^2}+\vartheta_u(b)\Bigr],
 \qquad \vartheta_u(b)=\frac{(u-2b)^4}{2b(u-b)u^2}\ge0,
\]
since $\frac1{2b(u-b)}=\frac2{u^2}+\frac{(u-2b)^2}{2b(u-b)u^2}$. By
Lemma~\ref{lem:A-AB}(a), the first part of $A_t$ is
\[
 \sum_{u=t+1}^{n-1}\frac{2n\,\E^{\br}_t(u-2M_u)^2}{(n-u)u^2}
 =\sum_{u=t+1}^{n-1}\frac{2n(t-1)^2(n-u)}{(n-t-1)^2u^2}
 +P_t\sum_{u=t+1}^{n-1}\frac{u-t-1}{u^2},
\]
with $P_t$ as in the proof of Lemma~\ref{lem:A-AB}(b). The terms of the
first sum are at most $8(t-1)^2/u^2$ and tend to $2(t-1)^2/u^2$, so the
first sum tends to $2(t-1)^2\psi'(t+1)$, where $\psi'(t+1)=\sum_{u>t}u^{-2}$.
The second sum is $H_{n-1}-H_t-(t+1)(\psi'(t+1)-\psi'(n))=\log n+\gamma-H_t-(t+1)\psi'(t+1)
+o(1)$, and $P_t=2+O(1/n)$. So the first part of
$A_t$ is $2\log n+2\gamma-2H_t+2((t-1)^2-t-1)\psi'(t+1)+o(1)$.

The second part is $\sum_{u=t+1}^{n-1}\frac n{n-u}\E^{\br}_t\vartheta_u(M_u)$.
For fixed $u$ the law of $M_u-1$ tends to $\mathrm{Bin}(u-t-1,\frac12)$ and
$\frac n{n-u}\to1$, so each term tends to $\E\vartheta_u(\bar M_u)$, where
$\bar M_u=1+N_u$ with $N_u\sim\mathrm{Bin}(u-t-1,\frac12)$. We dominate the
terms uniformly in $n$. We have $u-2M_u-\E(u-2M_u)=-2(Z-\E Z)$. The last
$n-u$ steps also form a uniform sample, so $h-M_u$ is hypergeometric with
$n-u$ draws and $u-2M_u-\E(u-2M_u)=2(h-M_u-\E(h-M_u))$. By Hoeffding's
comparison with sampling with replacement \cite[Theorem~4]{hoeffding1963},
applied to a convex fourth power, the fourth central moment of each is at
most that of the binomial law with the same mean, which is
$kpq(1+3(k-2)pq)\le k^2/2$ for $k$ draws. With $|\E(u-2M_u)|\le t$ and
$(x+y)^4\le8(x^4+y^4)$, this gives
$\E^{\br}_t(u-2M_u)^4\le64\min(u,n-u)^2+8t^4$. Also $\frac n{n-u}\le2u$,
because $u(n-u)\ge n-1\ge\frac n2$, and
$\frac n{n-u}\min(u,n-u)^2\le2u^2$, which is clear for $u\le\frac n2$ and
follows from $n(n-u)\le2u^2$ for $u>\frac n2$. On $\Omega_u$ we have
$2M_u(u-M_u)u^2\ge\frac7{32}u^4$, so the part of the $u$-th term on
$\Omega_u$ is at most $\frac{32}{7u^4}(128u^2+16t^4u)$. Off $\Omega_u$,
$M_u(u-M_u)\ge\frac u2$ gives $\vartheta_u\le u$, so that part is at most
$2u^2\Prob^{\br}_t(\Omega_u^c)$. This is at most $4u^2e^{-0.0943u}$ for
$u\ge8(t+1)$ by Lemma~\ref{lem:A-AB}(c), and at most $2u^2$ for the finitely
many smaller $u$. The bound is summable and does not depend on $n$. By
dominated convergence,
\[
 A_t-2\log n\to\tilde A_t:=2\gamma-2H_t+2\bigl((t-1)^2-t-1\bigr)\psi'(t+1)
 +\sum_{u>t}\E\vartheta_u(\bar M_u).
\]
Since $\E(u-2\bar M_u)^2=(t-1)^2+(u-t-1)$, the same computation with
$\frac n{n-u}$ replaced by 1 and $n$ by a cutoff $U$ shows that
$\tilde A_t=\lim_{U\to\infty}\bigl(\sum_{u=t+1}^U\bar\theta_u-2\log U\bigr)$,
where $\bar\theta_u=\E\frac{(u-2\bar M_u)^2}{2\bar M_u(u-\bar M_u)}$.

\emph{Part (b), the closed form.} Let $m=u-t-1$, $N=N_u$, $y=\frac{1+x}2$
and $\sigma(x)=\sum_{j=0}^{t-2}x^j$, with $\sigma=0$ for $t=1$. Since
$(u-2b)^2=u^2-4b(u-b)$ and $\frac1{b(u-b)}=\frac1u(\frac1b+\frac1{u-b})$,
we have $\bar\theta_u=\frac u2\E[\frac1{1+N}+\frac1{u-1-N}]-2$, and
$u-1-N$ has the law of $t+N$ because $N\overset d=m-N$. For $b\ge1$,
$\E\frac1{b+N}=\int_0^1x^{b-1}\E x^N\,dx=\int_0^1x^{b-1}y^m\,dx$. Hence
\[
 \bar\theta_u=u\int_0^1\frac{1+x^{t-1}}2\,y^m\,dx-2
 =\Bigl[\frac{2t}{m+1}-\Bigl(1+\frac t{m+1}\Bigr)2^{-m}\Bigr]-\mathcal T_m,
 \qquad\mathcal T_m=u\int_0^1(1-y)\sigma(x)y^m\,dx,
\]
using $\frac{1+x^{t-1}}2=1-(1-y)\sigma(x)$,
$\int_0^1y^m\,dx=\frac{2(1-2^{-m-1})}{m+1}$ and $u=m+1+t$. Split
$u=(m+1)+t$ in $\mathcal T_m$. Since $\frac d{dx}y^{m+1}=\frac{m+1}2y^m$,
integration by parts gives
$(m+1)\int_0^1x^j\frac{1-x}2y^m\,dx=-\delta_{j0}2^{-m-1}
+\int_0^1y^{m+1}((j+1)x^j-jx^{j-1})\,dx$ for $j\ge0$, and the integrands add
up to $(t-1)x^{t-2}$ over $0\le j\le t-2$. With
$I(j,k)=\int_0^1x^jy^k\,dx$ and $(1-y)\sigma(x)=\frac{1-x^{t-1}}2$,
\[
 \mathcal T_m=-[t\ge2]\,2^{-m-1}+(t-1)I(t-2,m+1)
 +\frac t2\int_0^1(1-x^{t-1})y^m\,dx .
\]
Sum over $0\le m\le m_0$, so that $U=m_0+t+1$. Since
$\sum_{m\le m_0}y^m=\frac{2(1-y^{m_0+1})}{1-x}$,
\begin{multline*}
 \sum_{u=t+1}^U\bar\theta_u=2tH_{m_0+1}
 -\sum_{m=0}^{m_0}\Bigl(1+\frac t{m+1}\Bigr)2^{-m}+[t\ge2](1-2^{-m_0-1})\\
 -(t-1)\sum_{k=0}^{m_0}I(t-2,k+1)-t\int_0^1\sigma(x)\bigl(1-y^{m_0+1}\bigr)dx .
\end{multline*}
Let $m_0\to\infty$. First $\sum_{m\ge0}(1+\frac t{m+1})2^{-m}=2+2t\log2$,
and $\int_0^1\sigma(1-y^{m_0+1})\,dx\to\int_0^1\sigma=H_{t-1}$. For
$t\ge2$ write $\sum_{k\le m_0}I(t-2,k+1)
=2\int_0^1x^{t-2}y\frac{1-y^{m_0+1}}{1-x}dx$ and
$x^{t-2}=1-(1-x^{t-2})$. The substitution $x=2y-1$ gives
$\int_0^1y\frac{1-y^{m_0+1}}{1-x}dx=\sum_{j=2}^{m_0+2}\frac{1-2^{-j}}j
=H_{m_0+2}-\frac12-\log2+o(1)$, and by dominated convergence
$\int_0^1(1-x^{t-2})y\frac{1-y^{m_0+1}}{1-x}dx\to
\int_0^1\sum_{j=0}^{t-3}x^j\frac{1+x}2dx=\frac12(H_{t-2}+H_{t-1}-1)$. So
$\sum_{k\le m_0}I(t-2,k+1)=2H_{m_0+2}-2\log2-H_{t-2}-H_{t-1}+o(1)$.
Collecting terms, with $2tH_{m_0+1}-2(t-1)H_{m_0+2}=2H_{m_0+1}+o(1)$,
\[
 \sum_{u=t+1}^U\bar\theta_u=2H_{m_0+1}-2-2\log2+[t\ge2]+(t-1)H_{t-2}-H_{t-1}+o(1).
\]
Here $[t\ge2]+(t-1)H_{t-2}-H_{t-1}=(t-2)H_{t-1}$ for all $t\ge1$ (for
$t=1$ the term $(t-1)H_{t-2}$ is absent), and
$H_{m_0+1}=\log U+\gamma+o(1)$. So
$\tilde A_t=2(\gamma-1-\log2)+(t-2)H_{t-1}=A^\infty_t$.

\emph{Part (c).} Since $\sum_tw_t=1$,
$k_2(n)-2\log n=\sum_{t\le h}w_t[A_t-2\log n-(t-1)]$. Both sides of
Proposition~\ref{prop:A-super}(c) equal $\pi_t$ at $\varepsilon=0$, and
$\kappa'(0)=2$. Comparing right derivatives at 0 and using part (a) gives
$A_t\le2(H_{n-1}-H_t)+\frac4{3t}+2(t+11)(\log\frac{(t+1)^2}{4t}+\frac23)$,
so $A_t-2\log n\le2\gamma+\frac43+2(t+11)\log(t+1)$ as in
Section~\ref{app:erw-proof-c}. Also $A_t-2\log n\ge-2\log t-6$ for
$t\le\frac{n+2}4$ by Lemma~\ref{lem:A-AB}(b), and $A_t\ge0$ gives
$A_t-2\log n\ge-2\log(4t)$ for $t>\frac{n+2}4$. With \eqref{eq:A-wt}, for
$n\ge64$ the $t$-th term is at most $C(1+t\log(t+1))t2^{-t}$ in absolute
value, with $C$ independent of $n$. By part (b),
Proposition~\ref{prop:first-mutation}(c) and dominated convergence,
$k_2(n)-2\log n\to\sum_{t\ge1}t2^{-t-1}[A^\infty_t-(t-1)]=K_2$.

To sum the series, note $\sum_tt2^{-t-1}=1$ and $\sum_tt(t-1)2^{-t-1}=2$.
Let $\mathcal H(x)=\sum_{k\ge0}H_kx^k=\frac{\ell}{1-x}$ with
$\ell=-\log(1-x)$, and $D=\frac d{dx}$. With $k=t-1$,
\[
 \sum_{t\ge1}t(t-2)2^{-t-1}H_{t-1}=\frac14\sum_{k\ge0}(k^2-1)H_k2^{-k}
 =\frac14\bigl[(xD)^2\mathcal H-\mathcal H\bigr]_{x=1/2}.
\]
Here $(xD)^2\mathcal H=x\bigl[\frac{1+\ell}{(1-x)^2}+\frac{x}{(1-x)^3}
+\frac{2x(1+\ell)}{(1-x)^3}\bigr]$ equals $8+6\log2$ at $x=\frac12$, and
$\mathcal H(\frac12)=2\log2$. So the sum is $2+\log2$, and
$K_2=2(\gamma-1-\log2)+(2+\log2)-2=2\gamma-2-\log2$. \qed

\subsection{Proof of the bound in \texorpdfstring{Remark~\ref{rem:refined}}{Remark 2.10}}\label{app:erw-refined}

Let $n\ge64$, $0<\varepsilon\le\frac12$ and $T=\lfloor n/8\rfloor$. We
repeat the lower bound of Section~\ref{app:erw-proof-c} but keep $A_t$. With
$\log(1-\varepsilon)\ge-\varepsilon-\varepsilon^2$,
$\sum_tw_t(t-1)\le2$ and $\sum_tw_t[A_t-(t-1)]=k_2(n)$, this gives
\[
 \log C_n\ge\log(\varepsilon c_1)+\varepsilon k_2(n)-\varepsilon\delta_n
 -\varepsilon^2(84\log n+1.05\cdot10^5),\qquad
 \delta_n=\sum_{t>T}w_tA_t .
\]
By the upper bound for $A_t$ in the proof of Theorem~\ref{thm:K2}(c) and
\eqref{eq:A-wt}, $\delta_n=O(n^2\log n\,2^{-n/8})$. Let
$F(x)=\log C_n(\frac xn)-\log E_n(\frac xn)$, and let $\hat F(x)$ be the
left side minus the right side of the equation for $\hat x_n$. Since
$\log(\varepsilon c_1)=\log\frac{4x}{n(n+2)}$ and
$\log E_n=\log\frac12+(n-1)\log(1-\frac xn)$, we get
$F(x)\ge\hat F(x)-\rho_n(x)$ for $0<x\le\frac n2$, with
$\rho_n(x)=\frac xn\delta_n+\frac{x^2}{n^2}(84\log n+1.05\cdot10^5)$. For
large $n$, $k_2(n)>0$ by Theorem~\ref{thm:K2}(c). Then $\hat F$ is strictly
increasing on $(0,n)$ with $\hat F'(x)=\frac1x+\frac{k_2(n)}n+\frac{n-1}{n-x}
\ge1$ on $[1,n)$, it tends to $+\infty$ as $x\to n$, and
$\hat F(1)=-2\log n+O(1)<0$. So $\hat x_n$ is its only zero, and
$\hat x_n>1$. By Proposition~\ref{prop:mlr}, $x_n<\frac n2$ and $F(x_n)=0$.
If $x_n>\hat x_n$, then
$x_n-\hat x_n\le\hat F(x_n)-\hat F(\hat x_n)\le\rho_n(x_n)$. Since
$x_n=O(\log n)$ by Theorem~\ref{thm:erw-crossing}, we get
$x_n\le\hat x_n+O((\log n)^3/n^2)$. \qed
